\documentclass[english,hidelinks,pdftex, 11 pt, class=report,crop=false]{standalone}
\usepackage[T1]{fontenc}
\usepackage[utf8]{luainputenc}
\usepackage{lmodern} % load a font with all the characters
\usepackage{geometry}
\geometry{verbose,a4paper, inner=0cm, outer=0 cm, bmargin=2cm, tmargin=1cm}
%\textwidth=12cm
\setlength{\parindent}{0bp}
\usepackage{import}
\usepackage[subpreambles=false]{standalone}
\usepackage{amsmath}
\usepackage{amssymb}
\usepackage{esint}
\usepackage{babel}
\usepackage{tabu}
\usepackage[dvipsnames, table]{xcolor}
\usepackage{cancel}
\makeatother
\makeatletter
\usepackage{datetime2}
\usepackage{titlesec}
\usepackage[many]{tcolorbox}

% Eheter
\newcommand{\enh}[1]{\,\textrm{#1}}
%referances
\newcommand{\net}[2]{{\color{blue}\href{#1}{#2}}}

%Spaces
\newcommand{\vsk}{\\[12pt]}
\newcommand{\vs}{\vspace{-12pt}}

% Tabell for opplegg

\newcommand{\ovlist}[1]{
\vspace{-16pt}
\begin{itemize}
	#1
\end{itemize}
}

% Chapters and sections
\titleformat{\section}[block]{\bfseries}{\hspace{3cm}\thesection}{5pt}{}
\titleformat{\subsection}[block]{\bfseries}{\hspace{3cm}\thesection}{5pt}{}
\newcommand{\sectionbreak}{\clearpage} % New page on each section
 

\newlength{\mywidth}
\setlength{\mywidth}{14cm}

\newcommand{\cont}[1]{
\begin{tcolorbox}[center, boxrule=0.0 mm, width=\mywidth,arc=0mm,enhanced jigsaw,,colback=white,breakable]
#1	
\end{tcolorbox}
}

\newcommand{\info}[5]{
\begin{tcolorbox}[center, boxrule=0.1 mm, width=\mywidth,arc=0mm,enhanced jigsaw,breakable,colback=yellow!5]	
	
	\footnotesize
	\textbf{Øvingsområde}\\[5pt] #1 
	
	\textbf{Utstyr}\\ #2  \\
	
	\begin{tabular}{@{} p{4cm} p{4cm} l} 
		\textbf{Tid} & \textbf{Elevinndeling} & \textbf{Læringsarena} \\
		#3  & #4 & #5
	\end{tabular} 
\end{tcolorbox}	
}

\newcommand{\gjen}[1]{\begin{tcolorbox}[center,boxrule=0.1 mm, width=\mywidth,arc=0mm,colback=blue!3] {\large \textbf{Gjennomføring} \vspace{5 pt}}\newline #1  \end{tcolorbox}\vspace{-5pt}}
\newcommand{\eks}[1]{\begin{tcolorbox}[center,boxrule=0.1 mm, width=\mywidth,arc=0mm,colback=green!3] {\large \textbf{Eksempel} \vspace{5 pt}}\newline #1  \end{tcolorbox}\vspace{-5pt}}

\newcounter{opl}
%\numberwithin{opl}{article}


\newcommand{\opl}[1]{
\newpage
{\refstepcounter{opl} %\phantomsection 
\large \textbf{\theopl \;#1} \vsk}
}

% Headlines
\newcommand{\fork}{\textbf{Forkunnskapar}\\}
\newcommand{\forb}{\textbf{Forberedelsar}\\}
\newcommand{\opgvr}{\textbf{Oppgaver}}



%colors
\newcommand{\colr}[1]{{\color{red} #1}}
\newcommand{\colb}[1]{{\color{blue} #1}}
\newcommand{\colo}[1]{{\color{orange} #1}}
\newcommand{\colc}[1]{{\color{cyan} #1}}
\definecolor{projectgreen}{cmyk}{100,0,100,0}
\newcommand{\colg}[1]{{\color{projectgreen} #1}}

% Lister med bokstavar
\usepackage[inline]{enumitem}
% Opg
\newcommand{\abc}[1]{
	\begin{enumerate}[label=\alph*),leftmargin=18pt]
		#1
	\end{enumerate}
}

\usepackage[]{hyperref}



\begin{document}
{\Large Eksamen P1 våren 2024\hfill {\footnotesize Løsning fra \color{blue} \href{https://sindreheggen.wordpress.com/}{OpenMathBooks prosjektet}}}	
\subsection*{Oppgave 1}
$ 2\text{\textperthousand}=0,002 $.
\alg{
2\text{\textperthousand\, av } x &= 18 \\
0.002\cdot x &= 18 \\
x &= \frac{18}{0.002} \\
&= 9\,000
}
9000 millioner er lik 9 milliarder. Befolkningen ville vært 9 milliarder.

\subsection*{Oppgave 2}
\abc{
\item 20\,000 er antall kroner Ada setter inn i banken. 1,0485 erd den årlige vekstfaktoren, som betyr at Ada får 4,85\% årlig rente. 
\item $ v $ i koden til Ada tilsvarer $ \frac{f(10)-f(0)}{10-0} $, som vil gi den gjennomsnittlige endringen til $ f $ på intervallet $ [0, 10] $. Dette vil fortelle hvor mange kroner sparepengene i snitt har økt med per år i løpet av 10 år. 
}

\subsection*{Oppgave 3}
\begin{itemize}
	\item For at to størrelser $ F $ og $ x $ skal være proporsjonale, må grafen til funksjonen $ F(x) $ være en rett linje som går gjennom origo. Dermed er $ f $ og $ x $ proporsjonale størrelser.
	\item Hvis $ F $ og $ x $ er omvendt proporsjonale størrelser, kan vi skrive $ F(x)=\frac{a}{x} $, hvor $ a $ er en konstant. Det betyr at $ F $ vil ha veldig høy verdi når $ x $ er i nærheten av 0. Når $ x $ har høy verdi, vil $ F $ nærme seg 0. $ F $ er ikke definert når $ x=0 $. Av dette ser vi at det bare er $ p $ og $ x $ som kan være omvendt proporsjonale størrelser.
\end{itemize}

\newpage
\subsection*{Oppgave 4}
\abc{
\item $ \frac{70}{10}=7 $. $ B=\frac{7^2}{2}=24,5 $.
\item \alg{
40,5 &= \frac{x^2}{2}\\
81 &= x^2 \\
\pm 9 &= x
}
$ x=9 $ tilsvarer en fart på $ 90\enh{km/h} $ ($ x $ ganget med 10).
}

\newpage
\subsection*{Oppgave 1}
\abc{
\item Ved regresjon finner vi $ f(x) $ som samsvarer med $O$ fra oppgaven.
\begin{figure}
	\includegraphics[scale=0.4]{opg1a}
\end{figure}
\item Finner ekstremalpunktet til $ f $. Overskuddet er størst ved ca. 283 solgte baguetter. Da er overskuddet ca. 4437 kr.
\begin{figure}
	\includegraphics[scale=0.3]{opg1b}
\end{figure}
\item Stigningstallet er ca 24. Dette forteller oss at mellom 100 og 200 solgte baguetter økte overskuddet med ca. 24 kr/baguette.
\begin{figure}
	\includegraphics[scale=0.4]{opg1c}
\end{figure}
}
\newpage
\subsection*{Oppgave 2}
\abc{
\item \textbf{Ved beregning} \\
$105607\cdot1.0325^5\approx 123\,920$ \vsk

\textbf{Grafisk} \\
Uttrykker beløpet som funksjonen $f$. Finner verdien til $f$ når $f$ skjærer $x=5$.
\begin{figure}
	\includegraphics[scale=0.4]{opg2a}
\end{figure}
Om fem år er det ca 123+,920 kr på kontoen.
\item \textbf{Ved beregning} \\
\alg{
x\cdot 1.0325^5&=105607 \\
 &= \frac{105607}{1.0325^5}\\
 &\approx 90\,000 
}
\textbf{Grafisk}
Fem år tilbake tilsvarer at $x=-5$. Finner verdien til $f$ når $f$ skjærer $x=-5$.
\begin{figure}
	\includegraphics[scale=0.4]{opg2b}
\end{figure}
For fem år siden satt Gaute inn 90\,000 kr.
}

\subsection*{Oppgave 3}
\abc{
\item $1,794\cdot10^6\cdot158,987=2.85222678\cdot10^8$. Det ble produsert ca. $2,85\cdot10^8$ liter olje i 2023.
\item $\frac{1,794}{1,685}\approx1,065$. Produksjonsmengden steg med ca. 6,5\% fra 2022 til 2023.
}

\subsection*{Oppgave 4}
\abc{
\item $3:27,14=3\cdot60+27,14\enh{s}=207,14\enh{s}$. Vi finner sekund per meter, og ganger med lengden av en runde:
\[ \frac{207,14}{1500\enh{m}}\cdot400\enh{m}\approx 55,24\enh{s} \]
Den gjennomsnittlige rundetiden var ca. 55,24\,s.
\item $ 25,89\enh{km/h}=25,89\cdot\frac{1000\enh{m}}{3600\enh{s}}\approx7,19\enh{m/s}$.\\
$3:43,73=223,73\enh{s}$.\\
$7,19\enh{m/s}\cdot223,73\enh{s}\approx 1609\enh{m}$.
1 engelsk mil er ca. 1690\,m.
}
\newpage
\subsection*{Oppgave 5}
\textbf{Alternativ 1}
\alg{
\text{Tall nr. 1} &= 2 \\
\text{Tall nr. 2} &= 5 = 2+3 \\
\text{Tall nr. 3} &= 11 = 5+6=5+3\cdot2 \\
\text{Tall nr. 4}&= 23 = 11+12=11+3\cdot2^2 \\
\text{Tall nr. 5} &= 47=23+24=23+3\cdot2^3 \\
\text{Tall nr. }n&= \text{Tall nr. }(n-1)+3\cdot2^{n-1} 
}
Tall nr. 2 er et oddetall. Alle tall etter tall nr. 2 kan uttrykkes som det forrige tallet addert med et produkt med 2 som faktor. Dette produktet er altså et partall. Et oddetall addert med et partall blir alltid et oddetall. \os

\textbf{Alternativ 2}
\alg{
\text{Tall nr. 1} &= 2 \\
\text{Tall nr. 2} &= 5 = 2+3=2+(2+1) \\
\text{Tall nr. 3} &= 11 = 5+6=5+(5+1)\\
\text{Tall nr. 4}&= 23 = 11+12=11+(11+1) \\
\text{Tall nr. 5} &= 47=23+24=23+(23+1) \\
\text{Tall nr. }n&= 2\cdot\text{Tall nr. }(n-1)+1
}
Leddet med faktor 2 vil alltid bli et partall. Når man så legger til 1 må summen bli et oddetall.
\newpage
\subsection*{Oppgave 6}
\abc{
	\item Skriver inn tabellen i rengearket i GeoGebra, lager liste med punkt, og bruker regresjon med potensfunksjon. Får da
	\[ K(x)=7.56x^{0.38} \]
	\item Aris modell:
	\[ f(x)=1000\cdot0.88^x \]
	der $ f $ er lufttrykket og $ x $ er km over havet. \os
	Lisas modell:
	\[ g(x)=1000\cdot0.5^\frac{x}{5.5} \]
	der $ g $ er lufttrykket og $ x $ er km over havet.
	\item Ved å skrive \texttt{Skjæring(K, y=85)}, finner vi at lufttrykket har verdi $ 578.34 $ når kokepunktet er $ 85^\circ $. Vi skriver så modellene til Ari og Lisa inn i GeoGebra, og finner at begge modellene gir en høyde på ca. 4.3 km over havet.
}
\localfig{opg6ac}{0.3}$1550\enh{kr}+1350\enh{kr}=1900\enh{kr}$
\localfig{opg6c}{0.3}

\subsection*{Oppgave 7}
Tilbudsprisen er $1550\enh{kr}+1350\enh{kr}=2900\enh{kr}$. Prisen uten tilbud er $1900\enh{kr}+800\enh{kr}=3700\enh{kr}$. 
Vi kan betale den samme prosentandelen av denne prisen som vi ville betalt av prisen uten tilbudet. Dette betyr at jeg betaler $\frac{800+1550}{3700}\approx 63,5\%$ av prisen. $2900\enh{kr}\cdot0,635=1841,5\enh{kr}$. Altså betaler jeg 1841,5\,kr, mens vennen min betaler 1058,5\,kr.

\end{document}